\documentclass[12pt]{article}
\usepackage[utf8]{inputenc}
\usepackage[T1]{fontenc}
\usepackage[margin=1in]{geometry}
\usepackage{setspace}
\usepackage{xcolor}
\usepackage[colorlinks=true, urlcolor=blue!60!black, linkcolor=blue!60!black]{hyperref}
\setlength{\parindent}{2em}
\setlength{\parskip}{1.25em}
\renewcommand{\baselinestretch}{1.0}
\begin{document}
\begin{center}
{\LARGE \textit{Hein v. Freedom From Religion Foundation, Inc.}, 551 U.S. 587 (2007)\par}
\vspace{0.8em}
{\large Case Brief\par}
\end{center}
\vspace{1em}

\section{Memory Jogger}

\textit{Hein} is the ``White House Faith-Based Conferences'' taxpayer-standing case (Course Unit: Standing; Case \#163). Federal taxpayers sued the directors of President George W. Bush's faith-based offices, alleging that taxpayer-funded executive conferences promoted religious social-service groups in violation of the Establishment Clause. In a fractured 5--4 decision, the Supreme Court held that the plaintiffs lacked Article~III standing because the narrow taxpayer-standing exception of \textit{Flast v. Cohen}, 392 U.S. 83 (1968), applies only to challenges against specific congressional exercises of the Taxing and Spending Clause, not discretionary Executive Branch expenditures paid out of general, lump-sum appropriations (\href{https://supreme.justia.com/cases/federal/us/551/587/}{551 U.S. 587}). Justice Alito's three-Justice plurality famously refused either to extend or to overrule \textit{Flast} (``We leave \textit{Flast} as we found it''), while Justice Scalia, joined by Justice Thomas, concurred in the judgment to blast the plurality's distinction as meaningless and urge overruling \textit{Flast} altogether (\href{https://supreme.justia.com/cases/federal/us/551/587/}{Justia}).

\section{Detailed Case Facts}

In 2001, President George W. Bush issued Executive Order No. 13199, 3 CFR 752 (2001 Comp.), creating the White House Office of Faith-Based and Community Initiatives within the Executive Office of the President (\href{https://supreme.justia.com/cases/federal/us/551/587/}{551 U.S. at 593--94}). The order stated that the office's mission was to ensure that private and charitable community groups, including religious ones, have the fullest opportunity permitted by law to compete on a level playing field for federal assistance while adhering to principles of pluralism, nondiscrimination, evenhandedness, and neutrality (\href{https://supreme.justia.com/cases/federal/us/551/587/}{Justia}). Through separate executive orders---including Executive Order Nos. 13198, 13279, 13280, 13342, and 13397---the President also established Executive Department Centers for Faith-Based and Community Initiatives inside several federal agencies and departments, instructing that direct federal assistance could not be used for inherently religious activities such as worship, religious instruction, or proselytization (\href{https://supreme.justia.com/cases/federal/us/551/587/}{Justia}).

Crucially, Congress enacted no statute specifically authorizing the creation of either the White House Office or the departmental centers, and Congress enacted no statute specifically appropriating funds for their activities (\href{https://supreme.justia.com/cases/federal/us/551/587/}{Justia}; \href{https://www.law.cornell.edu/supct/html/06-157.ZS.html}{Cornell LII}). Instead, the offices were funded entirely through general, lump-sum Executive Branch appropriations---for example, the Department of Education's center was funded from appropriations for the Office of the Secretary of Education, and the Department of Housing and Urban Development's center was funded from that Department's salaries and expenses account (\href{https://supreme.justia.com/cases/federal/us/551/587/}{Justia}).

The respondents were the Freedom From Religion Foundation, Inc. (FFRF), a Madison, Wisconsin nonstock corporation opposed to government endorsement of religion, and three of its members and officers---Anne Nicol Gaylor, Annie Laurie Gaylor, and Dan Barker (\href{https://ffrf.org/news/releases/7thcwin/}{FFRF}; \href{https://supreme.justia.com/cases/federal/us/551/587/}{Justia}). They sued Jay F. Hein (Director of the White House Office of Faith-Based and Community Initiatives, substituted for former Director Jim Towey) and the directors of various departmental centers in their official capacities (\href{https://supreme.justia.com/cases/federal/us/551/587/}{Justia}; \href{https://www.oyez.org/cases/2006/06-157}{Oyez}). Respondents alleged that petitioners violated the Establishment Clause by organizing national and regional conferences at which President Bush, former Secretary of Education Rod Paige, and other officials used religious imagery, praised the efficacy of faith-based programs, singled out faith-based groups as particularly worthy of federal funding, and sent a message to religious believers that they were favored political ``insiders'' while sending a message to nonbelievers that they were ``outsiders'' (\href{https://supreme.justia.com/cases/federal/us/551/587/}{Justia}). Respondents asserted no individualized economic, competitive, or direct-contact injury; their sole asserted basis for Article~III standing was that the individual respondents were federal taxpayers opposed to the Executive Branch's use of congressional appropriations to fund these conferences (\href{https://supreme.justia.com/cases/federal/us/551/587/}{Justia}).

\section{Procedural History}

\begin{itemize}
    \item \textbf{U.S. District Court for the Western District of Wisconsin:} In \textit{Freedom From Religion Foundation, Inc. v. Towey}, No. 04-C-381-S (W.D. Wis. Nov. 15, 2004), the District Court dismissed respondents' claims against the faith-based office directors for lack of Article~III standing (\href{https://supreme.justia.com/cases/federal/us/551/587/}{Justia}). The court concluded that under \textit{Flast v. Cohen}, 392 U.S. 83 (1968), federal taxpayer standing is strictly confined to challenges against exercises of congressional power under the Taxing and Spending Clause of Art.~I, \S~8, and that because petitioners acted on the President's behalf rather than administering a congressional spending program, plaintiffs lacked standing (\href{https://supreme.justia.com/cases/federal/us/551/587/}{Justia}; \href{https://www.oyez.org/cases/2006/06-157}{Oyez}).
    \item \textbf{U.S. Court of Appeals for the Seventh Circuit (Panel):} In \textit{Freedom From Religion Foundation, Inc. v. Chao}, 433 F.3d 989 (7th Cir. 2006), a divided panel reversed (\href{https://supreme.justia.com/cases/federal/us/551/587/}{Justia}). Writing for the majority, Judge Richard Posner held that \textit{Flast} permits federal taxpayers to challenge Executive Branch programs financed by general congressional appropriations whenever the marginal or incremental cost to the taxpaying public of the alleged Establishment Clause violation is greater than zero (\href{https://supreme.justia.com/cases/federal/us/551/587/}{Justia}; \href{https://ffrf.org/news/releases/7thcwin/}{FFRF}). Judge Kenneth Ripple dissented, warning that the majority's rule represented a dramatic expansion of taxpayer standing untethered from congressional taxing-and-spending enactments (\href{https://supreme.justia.com/cases/federal/us/551/587/}{Justia}).
    \item \textbf{Seventh Circuit (Rehearing En Banc Denied):} On May 3, 2006, the Seventh Circuit denied rehearing en banc by a 7--4 vote, \textit{Freedom From Religion Foundation, Inc. v. Chao}, 447 F.3d 988 (7th Cir. 2006) (\href{https://law.justia.com/cases/federal/appellate-courts/F3/447/988/505254/}{Justia 7th Cir.}). Concurring in the denial, Chief Judge Flaum and Judge Easterbrook expressed serious doubts about the coherence of Supreme Court taxpayer-standing precedent and invited the Supreme Court to resolve the tension, while Judge Ripple dissented from the denial of rehearing (\href{https://law.justia.com/cases/federal/appellate-courts/F3/447/988/505254/}{Justia 7th Cir.}).
    \item \textbf{U.S. Supreme Court:} The Court granted certiorari on December 1, 2006 (549 U.S. 1074), heard oral argument on February 28, 2007, and issued its decision reversing the Seventh Circuit on June 25, 2007 (\href{https://supreme.justia.com/cases/federal/us/551/587/}{Justia}; \href{https://www.oyez.org/cases/2006/06-157}{Oyez}).
\end{itemize}

\section{Judicial Votes}

The judgment reversing the Seventh Circuit was decided by a \textbf{5--4} vote, with no single opinion commanding a majority of the Court (\href{https://supreme.justia.com/cases/federal/us/551/587/}{Justia}; \href{https://www.law.cornell.edu/supct/html/06-157.ZS.html}{Cornell LII}):

\begin{itemize}
    \item \textbf{Plurality Opinion (3 Justices):} Justice Alito announced the judgment of the Court and delivered an opinion joined by Chief Justice Roberts and Justice Kennedy.
    \item \textbf{Concurring Opinion (1 Justice):} Justice Kennedy filed a separate concurring opinion (while also joining Justice Alito's plurality opinion in full).
    \item \textbf{Opinion Concurring in the Judgment (2 Justices):} Justice Scalia filed an opinion concurring in the judgment, joined by Justice Thomas (arguing that \textit{Flast} should be overruled).
    \item \textbf{Dissenting Opinion (4 Justices):} Justice Souter filed a dissenting opinion, joined by Justices Stevens, Ginsburg, and Breyer.
\end{itemize}

\section{Holding}

\textbf{Judgment:} Reversed. Respondents lacked Article~III standing as federal taxpayers to challenge the Executive Branch's Faith-Based and Community Initiatives conferences (\href{https://supreme.justia.com/cases/federal/us/551/587/}{551 U.S. at 615}).

\textbf{Rule Announced:} Under the narrowest-grounds principle for fragmented decisions (\textit{Marks v. United States}, 430 U.S. 188, 193 (1977)), Justice Alito's plurality opinion states the controlling rule: the narrow Establishment Clause exception to the general prohibition against federal taxpayer standing recognized in \textit{Flast v. Cohen}, 392 U.S. 83 (1968), does not extend to challenges against purely discretionary Executive Branch expenditures funded out of general, lump-sum appropriations (\href{https://supreme.justia.com/cases/federal/us/551/587/}{Justia}; \href{https://www.law.cornell.edu/constitution-conan/article-3/section-2/clause-1/taxpayer-standing}{Cornell LII ConAn}). To invoke \textit{Flast} standing, a federal taxpayer must challenge expenditures expressly authorized or mandated by a specific congressional enactment under the Taxing and Spending Clause of Art.~I, \S~8 (\href{https://supreme.justia.com/cases/federal/us/551/587/}{Justia}).

\textbf{Holding vs. Dicta:} The binding holding is confined to denying federal taxpayer standing where the challenged religious activity is funded through general executive administrative appropriations rather than a specific congressional spending mandate. Several important observations in the opinions constitute dicta rather than binding holdings: (1) Justice Alito's suggestion that Congress could quickly step in or that non-taxpayer plaintiffs might have standing to challenge hypothetical extreme executive abuses (such as building a federal house of worship); (2) Justice Alito's remark in Footnote~7 that informal legislative-history ``earmarks'' in committee reports do not establish legal requirements under \textit{Lincoln v. Vigil}, 508 U.S. 182 (1993), though lower courts treat this as strong guidance on statutory specificity; (3) Justice Kennedy's exhortation that executive and legislative officials bear an independent constitutional duty to obey the Establishment Clause even when no plaintiff has standing to sue; and (4) Justice Scalia's extensive critique advocating the overruling of \textit{Flast} and repudiation of ``Psychic Injury,'' which commanded only two votes (\href{https://supreme.justia.com/cases/federal/us/551/587/}{Justia}).

\section{Analysis of Opinions}

\subsection{Plurality Opinion (Justice Alito, joined by Chief Justice Roberts and Justice Kennedy)}

Justice Alito began from first principles of Article~III jurisdiction: federal courts are limited to actual ``Cases'' and ``Controversies,'' which require a plaintiff to show a concrete personal injury fairly traceable to the defendant's conduct and likely to be redressed by judicial relief (\href{https://supreme.justia.com/cases/federal/us/551/587/}{Justia}, citing \textit{Allen v. Wright}, 468 U.S. 737, 751 (1984)). Under \textit{Frothingham v. Mellon}, 262 U.S. 447 (1923), and \textit{Doremus v. Board of Education}, 342 U.S. 429 (1952), a federal taxpayer's interest in Treasury funds is shared with millions of others, minute, and indeterminable, making taxpayer status generally insufficient to confer Article~III standing (\href{https://supreme.justia.com/cases/federal/us/551/587/}{Justia}).

Turning to \textit{Flast v. Cohen}, 392 U.S. 83 (1968), Justice Alito emphasized that \textit{Flast}'s two-part nexus test---requiring a logical link between taxpayer status and a legislative enactment under the Taxing and Spending Clause of Art.~I, \S~8, and a nexus between taxpayer status and a specific constitutional limitation on that power (the Establishment Clause)---has been applied with ``rigor'' and largely confined to its facts (\href{https://supreme.justia.com/cases/federal/us/551/587/}{Justia}, citing \textit{Valley Forge Christian College v. Americans United for Separation of Church and State, Inc.}, 454 U.S. 464, 481 (1982)). In \textit{Flast}, Congress had expressly appropriated \$100 million under the Elementary and Secondary Education Act of 1965 and mandated that funds be provided on an equitable basis to private (predominantly religious) schools. Similarly, in \textit{Bowen v. Kendrick}, 487 U.S. 589 (1988), although the Court permitted an as-applied challenge to executive grant administration, the Adolescent Family Life Act expressly authorized grants and explicitly contemplated the involvement of religious organizations (\href{https://supreme.justia.com/cases/federal/us/551/587/}{Justia}).

By contrast, respondents in \textit{Hein} challenged no specific congressional statute or appropriation; they attacked discretionary executive spending drawn from general lump-sum accounts. Extending \textit{Flast} to purely executive expenditures would collapse its first prong because virtually every Executive Branch act---including speeches, proclamations, and daily operations---is funded by general congressional appropriations (\href{https://supreme.justia.com/cases/federal/us/551/587/}{Justia}). Such an expansion would violate separation-of-powers principles by enlisting federal courts as ``continuing monitors'' of the President and Cabinet officers. Justice Alito rejected both the Seventh Circuit's ``zero-marginal-cost'' test (noting Judge Easterbrook's point that presidential speeches and travel easily exceed \$500,000 annually) and respondents' ``fairly traceable'' expenditure test as unworkable (\href{https://supreme.justia.com/cases/federal/us/551/587/}{Justia}). Finally, invoking \textit{stare decisis} minimalism, the plurality declined to reconsider \textit{Flast}, concluding that a precedent need not be expanded to the outer limits of its logic (\href{https://supreme.justia.com/cases/federal/us/551/587/}{Justia}).

\subsection{Concurring Opinion (Justice Kennedy)}

Joining Justice Alito's opinion in full, Justice Kennedy wrote separately to affirm two distinct propositions: first, that \textit{Flast} was correctly decided and should not be overruled, because the Establishment Clause reflects a special constitutional concern that freedom of conscience not be compromised by government taxing and spending in aid of religion; and second, that separation-of-powers principles strictly forbid extending \textit{Flast} to day-to-day executive operations and public dialogues (\href{https://supreme.justia.com/cases/federal/us/551/587/}{Justia}). Permitting taxpayer suits over executive conferences and speeches would invite intrusive discovery and turn federal judges into ``speech editors for communications issued by executive officials and event planners for meetings they hold'' (\href{https://supreme.justia.com/cases/federal/us/551/587/}{Justia}). Justice Kennedy closed by reminding executive and legislative officials that lack of judicial standing does not excuse them from their independent oath-bound duty to obey the Constitution (\href{https://supreme.justia.com/cases/federal/us/551/587/}{Justia}).

\subsection{Opinion Concurring in the Judgment (Justice Scalia, joined by Justice Thomas)}

Justice Scalia delivered a blistering critique of both the plurality and the dissent, arguing that the Court's Establishment Clause taxpayer-standing jurisprudence was irredeemably incoherent because it oscillated between two incompatible theories of Article~III injury: ``Wallet Injury'' (higher personal tax liability, which fails traceability and redressability under \textit{Frothingham}) and ``Psychic Injury'' (mental displeasure that one's tax dollars are spent in aid of religion, which \textit{Flast} and \textit{Bowen} quietly embraced) (\href{https://supreme.justia.com/cases/federal/us/551/587/}{Justia}).

According to Justice Scalia, the Court faced a stark binary choice: either Psychic Injury satisfies Article~III---in which case \textit{Flast} logically covers all government expenditures alleged to violate the Establishment Clause, including executive spending---or Psychic Injury is an impermissible generalized grievance under \textit{Lujan v. Defenders of Wildlife}, 504 U.S. 555 (1992), in which case \textit{Flast} must be overruled (\href{https://supreme.justia.com/cases/federal/us/551/587/}{Justia}). He accused the plurality of inventing ``utterly meaningless distinctions'' divorced from Article~III's injury, traceability, and redressability criteria, noting that under the plurality's rule, taxpayers would lack standing even if the President diverted hospital grants to the Catholic Church in direct defiance of Congress, and that the plurality's rule flatly contradicted \textit{Bowen v. Kendrick} (\href{https://supreme.justia.com/cases/federal/us/551/587/}{Justia}). Concluding that \textit{Flast} rested on a repudiated view of standing that ignored separation of powers and generated no reliance interests, Justice Scalia argued that \textit{stare decisis} did not warrant ``beating \textit{Flast} to a pulp'' while leaving it technically alive; the Court should have overruled \textit{Flast} outright (\href{https://supreme.justia.com/cases/federal/us/551/587/}{Justia}).

\subsection{Dissenting Opinion (Justice Souter, joined by Justices Stevens, Ginsburg, and Breyer)}

Justice Souter argued that the respondents sought simply to apply \textit{Flast}, not extend it (\href{https://supreme.justia.com/cases/federal/us/551/587/}{Justia}). Rooted in James Madison's \textit{Memorial and Remonstrance Against Religious Assessments}---which warned against forcing a citizen to contribute even ``three pence'' to support a religious establishment---the Establishment Clause protects an individual right of conscience not to have identifiable tax funds extracted and spent to promote religion (\href{https://supreme.justia.com/cases/federal/us/551/587/}{Justia}). That injury is identical whether Congress mandates the spending or an executive agency directs it out of general appropriations.

Justice Souter argued that separation of powers offered no reason to insulate executive spending from Establishment Clause review while policing congressional spending, because courts owe equal respect to both coequal branches (\href{https://supreme.justia.com/cases/federal/us/551/587/}{Justia}). He contended that the plurality's distinction between ``congressional mandate'' and ``executive discretion'' was arbitrary and irreconcilable with \textit{Bowen v. Kendrick}, where the statute was upheld on its face and the as-applied challenge necessarily targeted executive discretion in selecting grantees (\href{https://supreme.justia.com/cases/federal/us/551/587/}{Justia}). Finally, analogizing to aesthetic harm (\textit{Friends of the Earth, Inc. v. Laidlaw Environmental Services}, 528 U.S. 167 (2000)), unequal bidding footing (\textit{Northeastern Florida Chapter of the Associated General Contractors v. Jacksonville}, 508 U.S. 656 (1993)), and racial gerrymandering (\textit{United States v. Hays}, 515 U.S. 737 (1995)), Justice Souter maintained that Article~III regularly recognizes intangible, noneconomic harms when tied to the specific constitutional interest protected (\href{https://supreme.justia.com/cases/federal/us/551/587/}{Justia}).

\section{Examples: Future Applications}

\begin{enumerate}
    \item \textbf{Same-Side (No Standing):} A federal agency uses funds from its general ``Salaries and Expenses'' lump-sum appropriation to pay for a nationwide series of workshops encouraging Christian charities to apply for federal disaster-relief contracts, featuring agency heads praising Christian theology. \textit{Analysis:} Like \textit{Hein}, there is no specific congressional enactment authorizing or mandating religious outreach, so federal taxpayers lack standing under \textit{Flast}.
    \item \textbf{Same-Side (No Standing):} The Department of Defense uses discretionary general operations funds to print and distribute evangelical devotional booklets to civilian visitors at Pentagon public tours, without any congressional statute directing or mentioning the program. \textit{Analysis:} Although the expenditure is identifiable and arguably violates the Establishment Clause on the merits, \textit{Hein} bars a federal taxpayer suit because the spending stems purely from executive discretion under a general appropriation (though a visitor directly subjected to unwelcome religious proselytization might pursue non-taxpayer standing).
    \item \textbf{Opposite-Side (Standing):} Congress enacts the ``Community Moral Renewal Act,'' specifically appropriating \$50 million to the Department of Health and Human Services and expressly directing that grants be awarded to ``houses of worship and sectarian ministries'' to provide youth counseling. Federal taxpayers sue, claiming the statute and grants violate the Establishment Clause. \textit{Analysis:} Unlike \textit{Hein}, this challenge targets a specific congressional appropriation and express statutory mandate under Art.~I, \S~8, squarely satisfying both prongs of \textit{Flast} and \textit{Bowen v. Kendrick}.
    \item \textbf{Opposite-Side (Standing):} A secular nonprofit social-service provider competes for a federal grant funded by general executive appropriations, is denied funding solely because the agency's faith-based office instructs reviewers to give preference to religious applicants, and sues under the Establishment Clause. \textit{Analysis:} The secular organization has standing---not as a taxpayer under \textit{Flast}, but under ordinary Article~III standing principles because it suffered a concrete, particularized competitive injury (loss of equal footing in the grant process), a pathway Justice Alito's plurality expressly left open (\href{https://supreme.justia.com/cases/federal/us/551/587/}{551 U.S. at 614}).
    \item \textbf{Fence-Sitter (Close Call):} Congress enacts a specific statutory grant program appropriating \$200 million for ``community rehabilitation organizations'' without mentioning religious groups in the statutory text, but the accompanying House Appropriations Committee report explicitly states that Congress expects at least \$40 million to be awarded to faith-based prison ministries, and the agency follows that guidance. Federal taxpayers bring an as-applied Establishment Clause challenge. \textit{Analysis:} This sits squarely on the razor's edge between \textit{Bowen v. Kendrick} and \textit{Hein}. Plaintiffs would argue that, as in \textit{Bowen}, Congress created a specific statutory grant-disbursement program under the Spending Clause rather than a general administrative account. The government would counter under \textit{Hein} Footnote~7 (\href{https://supreme.justia.com/cases/federal/us/551/587/}{Justia}) that nonstatutory committee-report earmarks have no legal force and that because the statute itself never mentioned or contemplated religious recipients, the choice to fund sectarian ministries reflects unreviewable executive discretion.
\end{enumerate}

\section{Critique}

\textbf{Originalist and Textualist Critique:} From an originalist and textualist perspective---articulated forcefully inside the Court by Justice Scalia and echoed by conservative scholars---the \textit{Hein} plurality committed a cardinal methodological error by refusing to overrule \textit{Flast}. Article~III vests ``The judicial Power'' only in ``Cases'' and ``Controversies,'' a limitation historically understood to exclude generalized public grievances about government spending (\href{https://supreme.justia.com/cases/federal/us/551/587/}{Justia}). Critics argue that by inventing an ``express congressional mandate'' boundary that has no logical connection to Article~III injury, traceability, or redressability, the plurality engaged in unprincipled ``stealth overruling'' or doctrinal ``narrowing'' (\href{https://columbialawreview.org/wp-content/uploads/2016/04/Re.pdf}{Re, 114 Colum. L. Rev. 1861}) that left lower courts trapped in what Judge Easterbrook called an ``arbitrary'' and ``illogical'' doctrine (\href{https://law.justia.com/cases/federal/appellate-courts/F3/447/988/505254/}{447 F.3d at 989--90}). As Justice Scalia later reiterated in \textit{Arizona Christian School Tuition Organization v. Winn}, 563 U.S. 125 (2011), \textit{Hein} distinguished \textit{Flast} away on unprincipled grounds instead of restoring Article~III's original meaning (\href{https://www.law.cornell.edu/supct/pdf/09-987P.ZC}{Cornell LII}).

\textbf{Progressive and Doctrinal Critique:} Progressive scholars and church-state separationists---including Ira Lupu and Robert Tuttle (\href{https://scholarship.law.gwu.edu/cgi/viewcontent.cgi?article=2151&context=faculty_publications}{GW Law}) and commentators in the \textit{Akron Law Review} (\href{https://ideaexchange.uakron.edu/akronlawreview/vol42/iss4/15/}{42 Akron L. Rev.})---criticize \textit{Hein} for creating a massive, formalistic loophole in Establishment Clause enforcement. Because modern governance relies heavily on broad lump-sum appropriations and executive agency discretion, \textit{Hein} effectively insulates vast swathes of federal religious promotion from judicial review whenever no discrete non-taxpayer plaintiff (such as a disappointed grant competitor) steps forward. Moreover, as Justice Souter and scholarly critics observe, the plurality's distinction between legislative and executive spending creates perverse incentives: Congress can shield religiously preferential programs from taxpayer suits simply by appropriating general funds and leaving religious allocation to the Executive Branch. In \textit{Winn}, 563 U.S. 125 (2011), the Court extended this restrictive trajectory by barring taxpayer standing to challenge state tuition tax credits, further shrinking \textit{Flast}'s reach (\href{https://www.law.cornell.edu/constitution-conan/article-3/section-2/clause-1/taxpayer-standing}{Cornell LII ConAn}; \href{https://www.law.cornell.edu/supct/html/09-987.ZD.html}{Kagan, J., dissenting}).

\textbf{Underlying Competing Values:} At bottom, \textit{Hein} reflects an enduring collision between two foundational constitutional commitments: (1) \textit{separation of powers and judicial restraint}, which warn that relaxing Article~III standing turns unelected judges into roving ombudsmen policing the daily speeches and administrative choices of democratically accountable officials; and (2) \textit{substantive enforcement of the Establishment Clause}, which warns that because an unlawful religious expenditure often harms the entire political community without inflicting a conventional pocketbook injury on any single citizen, denying taxpayer standing risks rendering a core First Amendment guarantee judicially unenforceable in the hands of the Executive Branch.

\section{Key Quotations}

\begin{quote}
``In light of the size of the federal budget, it is a complete fiction to argue that an unconstitutional federal expenditure causes an individual federal taxpayer any measurable economic harm. And if every federal taxpayer could sue to challenge any Government expenditure, the federal courts would cease to function as courts of law and would be cast in the role of general complaint bureaus.''
\end{quote}
\noindent --- Justice Alito (plurality opinion), \href{https://supreme.justia.com/cases/federal/us/551/587/}{551 U.S. at 593}.

\begin{quote}
``It is a necessary concomitant of the doctrine of \textit{stare decisis} that a precedent is not always expanded to the limit of its logic. That was the approach that then-Justice Rehnquist took in his opinion for the Court in \textit{Valley Forge}, and it is the approach we take here. We do not extend \textit{Flast}, but we also do not overrule it. We leave \textit{Flast} as we found it.''
\end{quote}
\noindent --- Justice Alito (plurality opinion), \href{https://supreme.justia.com/cases/federal/us/551/587/}{551 U.S. at 615}.

\begin{quote}
``It must be remembered that, even where parties have no standing to sue, members of the Legislative and Executive Branches are not excused from making constitutional determinations in the regular course of their duties. Government officials must make a conscious decision to obey the Constitution whether or not their acts can be challenged in a court of law and then must conform their actions to these principled determinations.''
\end{quote}
\noindent --- Justice Kennedy (concurring opinion), \href{https://supreme.justia.com/cases/federal/us/551/587/}{551 U.S. at 618}.

\begin{quote}
``If this Court is to decide cases by rule of law rather than show of hands, we must surrender to logic and choose sides: Either \textit{Flast v. Cohen}, 392 U. S. 83 (1968), should be applied to (at a minimum) all challenges to the governmental expenditure of general tax revenues in a manner alleged to violate a constitutional provision specifically limiting the taxing and spending power, or \textit{Flast} should be repudiated. For me, the choice is easy. \textit{Flast} is wholly irreconcilable with the Article III restrictions on federal-court jurisdiction that this Court has repeatedly confirmed are embodied in the doctrine of standing.''
\end{quote}
\noindent --- Justice Scalia (concurring in the judgment), \href{https://supreme.justia.com/cases/federal/us/551/587/}{551 U.S. at 618}.

\begin{quote}
``Here, the controlling, plurality opinion declares that \textit{Flast} does not apply, but a search of that opinion for a suggestion that these taxpayers have any less stake in the outcome than the taxpayers in \textit{Flast} will come up empty: the plurality makes no such finding, nor could it. Instead, the controlling opinion closes the door on these taxpayers because the Executive Branch, and not the Legislative Branch, caused their injury. I see no basis for this distinction in either logic or precedent, and respectfully dissent.''
\end{quote}
\noindent --- Justice Souter (dissenting opinion), \href{https://supreme.justia.com/cases/federal/us/551/587/}{551 U.S. at 637}.

\end{document}
